\subsection{Coherent changes of type and the sequential comparisons}

The distinction between an ind-isomorphism and a weak equivalence is material here. We use an enrichment of the actual indexing category to retain the coherence that an ind-object in an ordinary homotopy category discards. All path categories in C.P1--C.P4 have arbitrary interior vertical arrows.

\begin{result}{C.P1}{Proved}{The enriched category of types}
For types $\tau,\sigma$, let $P(\tau,\sigma)$ be the category of nondecreasing endpoint-preserving functors $I_\tau\to I_\sigma$, with natural transformations equal to the identity at the endpoints. This is a finite poset. When it is nonempty, its classifying space is contractible. Composition makes these categories the hom-categories of a strict $2$-category $\mathcal T$. Concatenation extends to a strict monoidal $2$-functor on $\mathcal T$.
\end{result}
\begin{proof}
Thinness follows from thinness of the target interval. We include the contraction from the companion's block at VII ms.~86 because its coherence, and not only its connectedness conclusion, is needed.

Induct on the number of constant-sign runs of $\sigma$. For an empty target, the mapping category is terminal. If the first source run has the opposite sign from the first target run, every map is constant on that source run; delete it, retaining its last vertex. Restriction is an isomorphism of mapping categories. We may thus assume that the two initial runs have the same sign $\varepsilon$. Write their vertex intervals as $[0,a]$ and $[0,b]$, with $a,b\geq1$. The vertex function $h(j)=\max(j,b)$ is an endofunctor of $I_\sigma$: it collapses the first run and fixes every subsequent edge.

For $F:I_\tau\to I_\sigma$, put $(RF)(0)=0$ and $(RF)(i)=hF(i)$ for $i>0$. A map cannot cross the first target turning vertex within the first source run. Thus $RF$ takes the first source run to $(0,b,\ldots,b)$, is a functor, remains nondecreasing, and preserves endpoints. Applying $h$ to the components away from zero defines $R$ on natural transformations. If $\varepsilon=+$ there is a natural transformation $1\Rightarrow R$; if $\varepsilon=-$ there is one $R\Rightarrow1$. All components lie in the first target run or are identities, so their common direction and naturality are valid. The image of $R$ is the full subcategory with prescribed first run $(0,b,\ldots,b)$, and $R$ is the identity on this subcategory. Deleting the prescribed initial runs identifies it with the corresponding mapping category for a target with one fewer run. Induction gives a finite natural-transformation contraction. This proves contractibility without claiming that the chosen contraction is strictly natural under every change of source type.

Vertical and horizontal composition of transformations preserve the endpoint and numerical-order conditions. Gluing at the common endpoint preserves the same conditions and the interchange law. These operations are the stated strict $2$-category and strict monoidal structure.
\end{proof}

\begin{result}{C.P2}{Proved}{Enriched descent and two final indexing functors}
Let $R$ be the reachability preorder of $\mathrm{Ty}$: its objects are types and it has one arrow $\tau\to\sigma$ exactly when $P(\tau,\sigma)$ is nonempty. For a small relative category $(M,\Sigma)$ with $\Sigma$ wide, put
\[
 F_\tau=\operatorname{Fun}_{\Sigma}(I_\tau,M)_{a,b}.
\]
The diagram $\tau\mapsto BF_\tau$ on $\mathrm{Ty}^{\mathrm{op}}$ admits a coherent extension along the functor
$\mathrm{Ty}^{\mathrm{op}}\to R^{\mathrm{op}}$. Moreover, this functor is homotopy final. If $t:\mathbb N^{\mathrm{op}}\to\mathrm{Ty}$ is a projective sequence such that every type receives a map from some $t_n$, its composite $\mathbb N\to R^{\mathrm{op}}$ is also homotopy final.
\end{result}
\begin{proof}
Precomposition is an actual functor
\[
 P(\tau,\sigma)\times F_\sigma\longrightarrow F_\tau,
 \qquad (u,c)\longmapsto c u.
\]
For a transformation $\alpha:u\Rightarrow v$ and a path transformation $\eta:c\Rightarrow d$, its component is the composite
$c u\xrightarrow{\eta u}d u\xrightarrow{d\alpha}d v$.
Its endpoint components are identities. Its interior components need not lie in $\Sigma$; this is the precise reason for using the full path categories. On objects, precomposition preserves the relative condition because a backward edge is sent to a string of backward edges or to an identity. The functors satisfy the enriched associativity and identity equations strictly.

Apply classifying spaces, using their natural compatibility with finite products. This gives a continuous enriched diagram on the topological category with hom-spaces $BP(\tau,\sigma)$. One completely specified model is to apply the singular functor to these hom-spaces and to $BF_\tau$, and to take the homotopy-coherent nerve of the simplicial category whose hom-spaces are $\operatorname{Sing}BP(\sigma,\tau)$ in the contravariant order. All hom-spaces are Kan complexes. C.P1 makes the map to $R^{\mathrm{op}}$ a Dwyer--Kan equivalence: it is the identity on objects and is a weak equivalence on every hom-space, including the empty ones. The coherent nerve therefore gives an equivalence of $\infty$-categories. Restriction along this equivalence is an equivalence of functor $\infty$-categories. Consequently the displayed enriched diagram is the pullback of a diagram $\overline F:R^{\mathrm{op}}\to\mathcal S$, up to a specified coherent equivalence. The space of such lifts with this equivalence retained is contractible. We use precisely these standard facts about coherent nerves and restriction, as stated in Kerodon, \S4.6.7, Theorem~4.8.4.1, and Proposition~4.5.4.8 \citep{ALurie2026}.

We verify finality using actual comma categories, rather than the reachability condition alone. Fix $\sigma\in R$. The comma category for $\mathrm{Ty}^{\mathrm{op}}\to R^{\mathrm{op}}$ at $\sigma$ is the opposite of the full subcategory
\[
 \mathrm{Ty}_\sigma=\{\tau\mid\operatorname{Hom}_{\mathrm{Ty}}(\tau,\sigma)\ne\varnothing\}.
\]
There is no choice of a map $\tau\to\sigma$ in this comma category: the target is the preorder $R$. The endofunctor $H(\tau)=\sigma*\tau$ of $\mathrm{Ty}_\sigma$ has natural transformations $H\Rightarrow1$ and $H\Rightarrow\operatorname{const}_\sigma$, obtained respectively by collapsing the first and second blocks. Their naturality follows directly from endpoint preservation and strict concatenation. They contract the nerve of $\mathrm{Ty}_\sigma$. Its opposite has the same homotopy type. The comma criterion proves finality.

For the sequence, the comma category at $\sigma$ is the full subposet of $\mathbb N$ consisting of $n$ for which $t_n\to\sigma$ exists. It is nonempty by hypothesis and upward closed because the transition maps are actual type maps. It therefore has a least element and contractible nerve. The same comma criterion proves finality. These are applications of the cofinality theorem with all comma hypotheses checked; see Kerodon, Theorem~7.2.3.1 \citep{ALurie2026}.
\end{proof}

\begin{result}{C.P3}{Proved}{All three comparisons of CATHOT E.26}
For every small $M$, every wide $\Sigma\subseteq M$, every $a,b\in M$, and every preorder-cofinal projective type sequence $t$, let $E_t$ be its fibred Grothendieck construction and $C_\infty(t)=\operatorname*{colim}_n F_{t_n}$. The specified functor
\[
 E_t\longrightarrow\operatorname{Ch}_\Sigma(M;a,b)
\]
is in $W_\infty$. The standard map $E_t\to C_\infty(t)$ is also in $W_\infty$. Thus its roof defines a canonical isomorphism
\[
 \beta_t:BC_\infty(t)\xrightarrow{\ \sim\ }
 B\operatorname{Ch}_\Sigma(M;a,b)\quad\text{in }\mathrm{Hot}.
\]
A strictly commuting map $u:t\to s$ of such projective sequences induces a weak equivalence $u^*:C_\infty(s)\to C_\infty(t)$, and
$\beta_t B(u^*)=\beta_s$. The comparisons obtained from $t*s\to t,s$ are independent of comparison choices and transitive in $\mathrm{Hot}$, with their actual maps and roofs understood as above. These assertions settle E.26(1), (2), and (3).
\end{result}
\begin{proof}
The natural fibred form of the Grothendieck homotopy-colimit theorem \citep[Corollary~26.9]{ChacholskiSchererPATHS} identifies
\[
 BE_t\simeq\operatorname*{hocolim}_{n\in\mathbb N}BF_{t_n},
 \qquad
 B\operatorname{Ch}_\Sigma(M;a,b)
 \simeq\operatorname*{hocolim}_{\tau\in\mathrm{Ty}^{\mathrm{op}}}BF_\tau.
\]
The maps in these comparisons are natural in both the indexing functor and the diagram. To check the variance, the opposite of the fibred construction is the ordinary covariant Grothendieck construction of the opposite fibre categories over the opposite base. The natural realization equivalence between a category and its opposite gives exactly the displayed formula. In particular, the map between these homotopy colimits is the image of the specified functor $E_t\to\operatorname{Ch}_\Sigma$, not an abstractly chosen equivalence.

Apply C.P2. Both homotopy colimits map equivalently to
$\operatorname*{colim}_{R^{\mathrm{op}}}\overline F$.
The comparison triangle commutes as a triangle of colimit cocones, with the coherent equivalence of diagrams in C.P2 included. Two-out-of-three proves that the specified map $E_t\to\operatorname{Ch}_\Sigma$ is a weak equivalence. The other map $E_t\to C_\infty(t)$ is a weak equivalence because nerves commute with filtered colimits of categories, and a sequential homotopy colimit of simplicial sets maps weakly equivalently to its ordinary filtered colimit. These arguments do not require fully faithful transition functors.

A natural transformation $u:t\to s$ produces the stated map of sequential colimits. The cartesian arrow
\[
 (t_n,F(u_n)c)\longrightarrow(s_n,c)
\]
in the all-type category gives a natural transformation between
$E_s\to E_t\to\operatorname{Ch}_\Sigma$ and the direct inclusion
$E_s\to\operatorname{Ch}_\Sigma$. Strict commutation of the sequence maps
is exactly its naturality in $n$. Naturality of the colimit cocones now
gives $\beta_t B(u^*)=\beta_s$. Since both $\beta$'s are isomorphisms,
$B(u^*)$ is one. The weak-equivalence assertion about the actual functor
follows from the definition of $W_\infty$.

Concatenating two projective sequences gives a projective sequence, with canonical block-collapse projections to both factors. It is preorder-cofinal, since a map from either factor to a type can be preceded by the corresponding block-collapse projection. Hence the resulting comparison is precisely $\beta_t^{-1}\beta_s$. All choices, after any required cofinal reindexing, give this same expression. Transitivity follows by cancelling the middle $\beta$. This proof retains the whole enriched action, so the known failure of an ind-isomorphism to imply an equivalence of unions does not apply.
\end{proof}

\begin{result}{C.P4}{Proved}{CATHOT E.27 for every admissible class and every cofinal block}
Let $\tau_0$ be a nonempty word containing both signs. Use the manuscript's left-padded, right-padded, and symmetrically padded systems of powers of $\tau_0$, with colimits $L_{\tau_0,\Sigma}$, $R_{\tau_0,\Sigma}$, and $P_{\tau_0,\Sigma}$. The actual maps $q_L,q_R$ from the first two colimits to the third are weak equivalences. On fixed endpoint fibres, the operation
\[
 \mu_{\tau_0}=B\kappa_{\tau_0}
       \bigl((Bq_L)^{-1}\times(Bq_R)^{-1}\bigr)
\]
is associative and unital in $\mathrm{Hot}$. It is independent of $\tau_0$ under the canonical comparisons of C.P3 and agrees with the strict all-type concatenation. No comparison with simplicial localization is needed for this conclusion.
\end{result}
\begin{proof}
The numbers of alternating runs of $\tau_0^n$ tend to infinity. Every finite sign word is obtained from a sufficiently long power by collapsing unused edges and composing or expanding retained same-sign runs. Thus all three projective sequences satisfy C.P3. The constant-padding maps arise contravariantly from strictly commuting maps of these sequences: for $q_L$, collapse the rightmost $n$ blocks of $\tau_0^{2n}$, and for $q_R$, collapse the leftmost $n$ blocks. They commute with the specified outer-collapse transition maps block by block. C.P3 makes $q_L,q_R$ weak equivalences.

Write $A_{a,b}=B\operatorname{Ch}_\Sigma(M;a,b)$ and let $\gamma_{a,b,c}:A_{a,b}\times A_{b,c}\to A_{a,c}$ be strict concatenation after realization, with factors written in traversal order: $\gamma(c,d)$ first traverses $c$ and then $d$. Thus, if enriched composition is written with factors $A_{b,c}\times A_{a,b}$ and denoted by $*$, the present $\gamma(c,d)$ is $d*c$. The comparisons of C.P3 satisfy
\[
 \beta_P Bq_L=\beta_L,\qquad
 \beta_P Bq_R=\beta_R,
 \qquad
 \beta_P B\kappa_{\tau_0}
       =\gamma(\beta_L\times\beta_R).
\]
For the last identity, evaluate on the common-stage diagram with two words $\tau_0^n$. Both maps send the pair to its literal concatenation, of type $\tau_0^{2n}$. Increasing $n$ pads the two outer ends, so the maps commute with the stated transitions. The diagonal $\mathbb N\to\mathbb N\times\mathbb N$ is final, and finite products of spaces commute with filtered colimits. Consequently this common-stage verification identifies the actual maps on the two sequential homotopy colimits. The natural Grothendieck comparison then gives the displayed equality in $\mathrm{Hot}$.

Substituting the first two identities into the third gives
\[
 \beta_P\mu_{\tau_0}
   =\gamma(\beta_P\times\beta_P).
\]
Strict associativity and the empty-path units for $\gamma$ imply associativity and the corresponding unit identities for $\mu_{\tau_0}$. Every other cofinal block gives the same transported operation on $A_{a,b}$. The construction uses all the coherent comparison maps of C.P2--C.P3, rather than an inference from equality on component-valued finite tests. Words containing only one sign have bounded useful length and are outside the statement; they are not the cofinal blocks used at VII ms.~125--127.
\end{proof}

\subsection{A finite obstruction in \texorpdfstring{$\mathrm{Cat}\,\mathrm{hot}$}{Cat hot}}

\begin{result}{C.P5}{Refuted}{The one-sided comparison need not be invertible modulo natural transformations}
There is a four-object category $X$, with $\Sigma=\operatorname{Mor}(X)$ and fixed endpoints $a=b$, for which the actual constant-tail extension
$q_L:\mathcal L(X;a,a)\to\mathcal P(X;a,a)$ of CATHOT E.23 is a weak equivalence but is not invertible in $\mathrm{Cat}\,\mathrm{hot}$. Thus the sequential comparisons cannot in general be promoted from $\mathrm{Hot}$ to $\mathrm{Cat}\,\mathrm{hot}$, even for finite $X$ and the canonical two-letter block.
\end{result}
\begin{proof}
Take $X$ to be the poset on $\mathbb Z/4$ whose only nonidentity arrows are
\[
 0\to1,\quad0\to3,\quad2\to1,\quad2\to3.
\]
There are four objects and eight arrows; compositions involving a nonidentity arrow have an identity as the other factor. Its nerve is the four-cycle, so $BX\simeq S^1$. Let $\widetilde X$ be the integer fence with arrows from an even integer to its two adjacent odd integers. Reduction modulo four is a covering functor $\widetilde X\to X$: every incoming and outgoing base arrow has a unique lift at a chosen vertex. Its nerve is the universal covering line.

Use the same integer fence as the parameter category $J$. Objects of $\mathcal P(X;0,0)$ are functors $J\to X$ eventually constant at zero on both tails; arrows are endpoint-fixed natural transformations. Every such path lifts uniquely to $\widetilde X$ with left tail zero. A natural transformation lifts uniquely with the same left-tail component, by the arrow-lifting property and its commuting squares. Its right-tail component is an identity, so the two lifted right endpoints agree. The lifted right endpoint is therefore constant on each zigzag component. Denote by $\mathcal P_0$ the full subcategory on paths whose right endpoint lift is zero. This definition, rather than a chosen identification of component sets with $\mathbb Z$, is all that is needed.

Suppose a functor $r:\mathcal P(X;0,0)\to\mathcal L(X;0,0)$ were an inverse to $q_L$ in $\mathrm{Cat}\,\mathrm{hot}$. In particular, there would be a finite zigzag of natural transformations, of some length $k$, between $1_{\mathcal P}$ and $q_Lr$. Evaluating this zigzag at a path $c\in\mathcal P_0$ produces a zigzag of $k$ endpoint-fixed path transformations. Every intermediate path also has right endpoint lift zero, since each transformation preserves that endpoint. Lift this objectwise zigzag to $\widetilde X$ and then evaluate it at any integer $i\geq0$. The final path $q_Lr(c)$ is constant at zero on the nonnegative half-line; because its right endpoint lift is zero, its lift there is also zero. Each arrow of $\widetilde X$, in either direction of a zigzag, changes the integer vertex by at most one. Hence the lifted value of $c$ at every $i\geq0$ would have distance at most $k$ from zero.

For $N\geq1$, define a path in $\widetilde X$ by
\[
 \widetilde c_N(i)=
 \begin{cases}
 0,&i\leq0,\\
 i,&0\leq i\leq4N,\\
 8N-i,&4N\leq i\leq8N,\\
 0,&i\geq8N.
 \end{cases}
\]
At every generating edge of $J$ this is an arrow of the integer fence or an identity. Its reduction modulo four is an object $c_N\in\mathcal P_0$, and its lifted value at the positive coordinate $4N$ is $4N$. Choosing $4N>k$ contradicts the preceding bound. No functor $r$ can therefore give the required one-sided inverse in $\mathrm{Cat}\,\mathrm{hot}$.

The weak equivalence of $q_L$ is C.P4 (and is also the independently checked CATHOT E.25). The program \texttt{code/b\_crown\_checks.py} checks the full finite base category, the covering rule at all residues, all endpoint-fixed transformations of parameter length six, and the explicit path family through $N=8$. The distance argument proves the obstruction for arbitrary $N$ and arbitrary candidate inverse functors; finite enumeration is not used as a substitute for that quantifier.
\end{proof}
